\documentclass[conference]{IEEEtran}
\IEEEoverridecommandlockouts
\usepackage{cite}
\usepackage{amsmath,amssymb}
\usepackage{graphicx}
\usepackage{booktabs}
\usepackage{textcomp}
\usepackage{xcolor}
\usepackage{url}
\def\UrlFont{\footnotesize\ttfamily}

% Placeholder for numbers that come from running the pipeline on the real
% data. Search the file for \tbd before submitting; none should remain.
\newcommand{\tbd}[1]{\textcolor{red}{[#1]}}

\begin{document}

\title{Milestone 1: Federated Intrusion Detection Under Label-Flipping Poisoning on CICIoT2023}

\author{
\IEEEauthorblockN{Salvador Leon Lopez, Elena Victoria Sanchez, Joe Reyna, April Flores (Group 8)}
\IEEEauthorblockA{Department of Computer Science, The University of Texas Rio Grande Valley\\
CSCI 6330: Foundations of Intelligent Security Systems, Fall 2026}
}

\maketitle

\section{Research Question}
Organizations that each see different attack traffic could train a shared
intrusion detector with federated learning (FL) without pooling raw captures
\cite{mcmahan2017}, but the server cannot inspect client data, so a
compromised participant can poison the shared model \cite{tolpegin2020}.

\textbf{RQ:} In a ten-client federated IDS trained on non-IID partitions of
CICIoT2023, how much global detection performance is lost when one client
performs a label-flipping attack, and does coordinate-wise trimmed-mean
aggregation \cite{yin2018} recover it?

\textbf{Expectation (falsifiable):} Under FedAvg, global macro F1 falls by
more than 3 points versus the paired clean run, and trimmed mean recovers at
least half of that loss. Smaller effects falsify it. Because label flipping
need not produce coordinate-wise extreme updates, weak recovery would itself
be a meaningful result. A secondary question varies the malicious share
(1, 2 and 3 of 10 clients).

\section{Dataset Selection}

\subsection{Why CICIoT2023}
CICIoT2023 \cite{neto2023} provides labeled flows with 46 numeric features
from 33 attacks on a 105-device IoT testbed, grouped by its authors into seven
families. It is recent, multi-class and heavily imbalanced, which makes macro
F1 and per-family miss rates informative. Its page states no license and asks
for citation; we use it for coursework only (accessed Sep.~20, 2026). It is
single-testbed lab data, so our clients are \emph{simulated} and all
heterogeneity is \emph{simulated statistical heterogeneity}, not evidence
about real organizations.

\subsection{Label Families}
Table~\ref{tab:families} gives our mapping, which follows the authors'
taxonomy; a unit test checks the family sizes and matching is
case-insensitive.

\begin{table}[t]
\caption{Label mapping and subset counts}
\label{tab:families}
\centering\footnotesize
\setlength{\tabcolsep}{3pt}
\begin{tabular}{@{}lp{3.3cm}rr@{}}
\toprule
Family & Original labels (count) & Raw rows & Subset \\
\midrule
Benign & BenignTraffic (1) & \tbd{X} & \tbd{X} \\
DDoS & 12 floods and fragmentations & \tbd{X} & \tbd{X} \\
DoS & UDP, SYN, TCP, HTTP floods (4) & \tbd{X} & \tbd{X} \\
Recon & 4 Recon-* + VulnerabilityScan (5) & \tbd{X} & \tbd{X} \\
Web & SQLi, XSS, CmdInj, Upload, Backdoor, BrowserHijack (6) & \tbd{X} & \tbd{X} \\
BruteForce & DictionaryBruteForce (1) & \tbd{X} & \tbd{X} \\
Spoofing & DNS, MITM-ARP (2) & \tbd{X} & \tbd{X} \\
Mirai & greeth, greip, udpplain (3) & \tbd{X} & \tbd{X} \\
\bottomrule
\end{tabular}
\end{table}

\subsection{Subset and Split Procedure}
The order is fixed and scripted (\texttt{prepare\_data.py}):
(1)~read all \tbd{N} CSV files and record label counts \emph{before} sampling;
(2)~sample at most $\lceil 2\cdot45{,}000/N\rceil$ rows per family per file
so every capture file contributes;
(3)~drop rows with inf/NaN (\tbd{n}), exact feature+label duplicates
(\tbd{n}) and feature vectors with conflicting labels (\tbd{n});
(4)~cap each family at 45{,}000 rows, keeping rarer families whole
(\tbd{total} rows). The subset is therefore closer to balanced than the raw
data; we report both prevalences and never present the subset's as natural;
(5)~hold out 20\% test and 10\% validation, stratified by family and
\emph{grouped by source file}, so no file contributes rows to two splits
(the CSVs expose no device or session field, so the file is the finest
capture grouping available). Only then are the scaler and partitioners fit,
on the 70\% training pool. Output hashes are stored in a manifest so every
run can be traced to the exact data.

\subsection{Client Partitions}
Following the feedback, we use ten clients and a controlled \textbf{Dirichlet}
label-skew split as the primary non-IID setting \cite{hsu2019}: for each
family, client shares are drawn from $\mathrm{Dir}(\alpha{=}0.5)$, redrawn
until every client holds $\geq$2\% of the pool. The label-free k-means
partition ($k{=}10$) from the proposal is retained as a reported alternative,
and an IID split is the control. Validation and test rows are assigned to
clients with the \emph{frozen} train-time partitioner (Dirichlet proportions,
or k-means scaler and centroids), never refit, so per-client held-out
evaluation is possible. For every client we report size, benign share,
family counts, label purity, and total-variation distance from the global
distribution, and check: share $\geq$2\%, size ratio $\leq$20, purity
$\leq$95\%, and $\geq$2 families with $\geq$50 rows (Table~\ref{tab:clients}).

\begin{table}[t]
\caption{Partition checks (training pool, 10 clients)}
\label{tab:clients}
\centering\footnotesize
\begin{tabular}{@{}lccc@{}}
\toprule
 & Dirichlet & k-means & IID \\
\midrule
Smallest / largest client (rows) & \tbd{X} & \tbd{X} & \tbd{X} \\
Max label purity & \tbd{X} & \tbd{X} & \tbd{X} \\
Mean TV distance from global & \tbd{X} & \tbd{X} & \tbd{X} \\
Normalized client--label MI & \tbd{X} & \tbd{X} & \tbd{X} \\
Checks & \tbd{PASS/FAIL} & \tbd{X} & \tbd{X} \\
\bottomrule
\end{tabular}
\end{table}

\section{Tool Choices and Method}

\subsection{Tools}
Python 3.11--3.13, PyTorch, scikit-learn and pandas, with pinned versions.
The FL loop is a single-process simulation written by the team
(\texttt{run\_federated.py}), which the project specification permits,
instead of Flower \cite{beutel2020}: Flower~1.39 deprecated its
\texttt{run\_simulation()} API, and a short loop makes every aggregation
step inspectable. FedAvg, trimmed mean and median are unit-tested.
\tbd{Hardware: CPU/GPU, RAM}. All runtimes are single-machine simulation time.

\subsection{Model and Training}
One MLP is held constant across all conditions: 46 inputs (signed log
transform, then standardization fit on the training pool), hidden layers of
128 and 64 units with ReLU and dropout 0.1, and 8 output logits (14{,}792
parameters). Loss is unweighted cross-entropy; Adam, learning rate
$10^{-3}$, batch size 256. The centralized reference trains 30 epochs; FL
runs 30 rounds of one local epoch with all ten clients each round and a fresh
optimizer state per round. The final epoch/round model is evaluated.

\subsection{Threat Model and Defense}
The attacker is the client with the most attack rows in the frozen training
partition (identical for every seed and condition) and must hold $\geq$2
attack families with $\geq$100 rows each. It relabels 50\% of its attack rows
\emph{within each family} as benign; the poisoned set is drawn once per seed
and reused every round. It changes labels only and otherwise behaves honestly.
Trimmed mean uses $\beta{=}0.2$: per coordinate, the $\lfloor 0.2\cdot10\rfloor{=}2$
highest and 2 lowest client values are removed and the remaining 6 averaged
without size weighting. It is also run without attack, since heterogeneous
honest clients can look like outliers.

\subsection{Conditions, Seeds and Metrics}
(1)~Centralized; (2a/2b)~clean FedAvg, IID and non-IID; (3)~FedAvg with one
attacker; (4a/4b)~trimmed mean, clean and attacked; plus 2 and 3 attackers.
The data seed (42) fixes records, splits and partitions for all runs. Each run
seed (42--44) changes only model initialization, minibatch order and the
poisoned draw, so conditions with the same seed start from identical weights
and are compared as paired differences; we report individual runs as well as
mean and standard deviation. The primary metric is global test macro F1, with
per-family precision, recall, FNR and FPR (each with support), binary
benign-vs-attack FPR/FNR, PR-AUC, per-client held-out macro F1, and
communication: 59.2\,kB per client per round in each direction, 1.18\,MB per
round and 35.5\,MB per run (raw float32, downloads included).

\section{Response to Proposal Feedback}
\begin{table}[t]
\caption{Feedback and changes}
\label{tab:feedback}
\centering\footnotesize
\setlength{\tabcolsep}{3pt}
\begin{tabular}{@{}p{2.6cm}p{5.5cm}@{}}
\toprule
Feedback & Change \\
\midrule
k-means clusters may be unequal or label-pure & Dirichlet split is primary; per-client reports and checks for every partition \\
Use 10--20 clients; vary malicious share & 10 clients; 1, 2, 3 attackers \\
Dedup and grouping before split; freeze partitioner & Duplicates removed first; file-grouped split; frozen held-out assignment \\
Global and client-level evaluation with support & Per-client held-out metrics; support on every rate \\
Deterministic attack & Fixed attacker, fixed per-seed poisoned set, labels only \\
Specify trimmed mean & $\beta{=}0.2$: 2 dropped per side, unweighted mean of 6 \\
Paired factors and seeds & Fixed data seed; identical initial weights per seed; paired deltas \\
Pilot, mapping, counts, hyperparameters & Pilot configuration; Table~\ref{tab:families}; Section~III-B \\
Communication and hardware & Bytes per client, per round and per run; hardware logged \\
\bottomrule
\end{tabular}
\end{table}
Table~\ref{tab:feedback} summarizes how each point of the feedback was
addressed. The largest design change is the move from six k-means clients to
ten Dirichlet clients; k-means remains as an alternative so the effect of the
partition can still be discussed.

\section{Status and Next Steps}
\textbf{Done:} the full pipeline (subset, splits, three partitioners with
checks, centralized model, FL loop, attack, FedAvg, trimmed mean, median,
summary tables) is implemented, and 22 unit and end-to-end tests pass.
\tbd{Pilot (12 files, 5 rounds): centralized validation macro F1 = X;
FedAvg IID = X, non-IID = X.} \tbd{Full subset: centralized validation macro
F1 = X; lowest per-family F1: family X.}
\textbf{Next (Milestone 2, Nov.~15):} run clean FedAvg on the full data with
three seeds, then conditions 3, 4a and 4b and the attacker sweep, reporting
paired differences. The test set is used once, after all settings are frozen.
\textbf{Limitations:} single-testbed lab data; simulated clients; a capped
subset whose prevalence differs from the raw data; three seeds support
paired comparisons, not distributional claims.

\footnotesize
\begin{thebibliography}{00}
\itemsep 0pt
\bibitem{mcmahan2017} B. McMahan \emph{et al.}, ``Communication-efficient learning of deep networks from decentralized data,'' \emph{Proc. AISTATS}, 2017.
\bibitem{tolpegin2020} V. Tolpegin, S. Truex, M. E. Gursoy, and L. Liu, ``Data poisoning attacks against federated learning systems,'' \emph{Proc. ESORICS}, 2020.
\bibitem{yin2018} D. Yin, Y. Chen, K. Ramchandran, and P. Bartlett, ``Byzantine-robust distributed learning,'' \emph{Proc. ICML}, 2018.
\bibitem{neto2023} E. C. P. Neto \emph{et al.}, ``CICIoT2023: A real-time dataset and benchmark for large-scale attacks in IoT environment,'' \emph{Sensors}, vol.~23, no.~13, 2023.
\bibitem{hsu2019} T.-M. H. Hsu, H. Qi, and M. Brown, ``Measuring the effects of non-identical data distribution for federated visual classification,'' arXiv:1909.06335, 2019.
\bibitem{beutel2020} D. J. Beutel \emph{et al.}, ``Flower: A friendly federated learning framework,'' arXiv:2007.14390, 2020.
\end{thebibliography}

\end{document}
